\begin{abstract}

    \begin{center}
        \large{% Единое место для указания названия работы
Предобработка мешей для использования в ускоряющих структурах
} \\
    \large\textit{Мусько Денис Владимирович} \\[1 cm]
    \end{center}

    Работа посвящена ускорению трассировки лучей в реальном времени за счёт
    предобработки геометрии перед подачей в структуру нижнего уровня (BLAS).
    Графический конвейер часто выделяет отдельный BLAS под каждый материал;
    если треугольники материала разбросаны по крупной модели, его AABB
    охватывает всю модель, а в TLAS возникают огромные перекрывающиеся
    инстансы, и луч вынужден пересекать почти все BLAS сцены.

    Цель --- разработать и оценить метод рантайм-предобработки, повышающий
    качество BLAS без визуальных потерь. Предложена глобальная
    пространственная кластеризация: треугольники всех материалов объединяются,
    сортируются по Morton-коду центроида и разбиваются на K компактных
    кластеров, каждый из которых регистрируется как отдельный BLAS с малым
    AABB. Метод работает на стороне CPU, не требует изменений драйвера и
    сохраняет попиксельную идентичность изображения.

    Метод реализован в семпле \texttt{testGI} движка DagorEngine. На «плохом»
    разбиении модели время \texttt{rtsm::do\_trace} падает в 8,6 раз
    (с 9,15 до 1,07~мс), на реалистичном --- в 2,7 раза
    (с 2,75 до 1,01~мс). Корректность подтверждена сравнением картин
    Acceleration Structure в Nsight Graphics. Рекомендуемый диапазон ---
    K от 64 до 128.


\end{abstract}
